\begin{solapartat}
L'energia potencial d'un moviment vibratori harmònic és $E_p = \frac{1}{2}kx^2$ \punts{0,1}, per tant:
\[ E_p = \frac{1}{2}kx^2 \Rightarrow 4 = \frac{1}{2}\,k\,1^2 \Rightarrow k = 8,00\ \mathrm{N/m}\ \text{\punts{0,2}} \]
Per altre banda tindrem:
\[ E_c = \frac{1}{2}mv^2 \Rightarrow 12 = \frac{1}{2}m(5,44)^2 \Rightarrow m = \frac{24}{5,44^2} = 8,11\cdot 10^{-1}\,\mathrm{kg}\ \text{\punts{0,2}} \]
Per l'energia total tinrem:
\[ E = E_p + E_c = 12 + 4 = 16\,\mathrm{J}\ \text{\punts{0,5}} \]
\end{solapartat}

\begin{solapartat}
La freqüència angular del moviment és $\omega^2 = k/m$ \punts{0,1} i per tant
\[ \omega = \sqrt{\frac{k}{m}} = \sqrt{\frac{8}{0,811}} = 3,14\,\mathrm{rad/s}\ \text{\punts{0,1}} \]
L'amplitud surt de l'expressió de la energia total del moviment:
\[ E = \frac{1}{2}kA^2\ \text{\punts{0,1}} \Rightarrow A = \sqrt{\frac{2E}{k}} = \sqrt{\frac{2\cdot 16}{8}} = 2,00\,\mathrm{m}\ \text{\punts{0,1}} \]
Per trobar la fase inicial hem d'anar a les condicions inicials, tot tenint present que l'equació general del moviment harmònic simple és
\[ x(t) = A\,\cos(\omega\,t + \varphi)\ \text{\punts{0,1}} \]
(També considerem correcte les expressions si partim de la elongació amb la funció sinus) i, per tant, a $t = 0$ resulta $x = A\cos\varphi$, $v = -A\omega\sin\varphi$. Amb $x(0) = 1\ \mathrm{m}$ i $v(0) = -5,44\ \mathrm{m/s}$, i els valor anteriors, obtenim
\[ 1 = 2\cos\varphi\ ;\ -5,44 = -2\cdot 3,14\sin\varphi \Rightarrow \cos\varphi = 0,5\ ;\sin\varphi = 0,86 \Rightarrow \tan\varphi = \frac{0,86}{0,5}\ \text{\punts{0,1}} \Rightarrow \]
\[ \varphi = \arctan\left(\frac{0,86}{0,5}\right) = 1,04\ \mathrm{rad}\ \text{\punts{0,1}} \]
Per tant l'equació del moviment és:
\[ x(t) = 2\cos(3,14\ \mathrm{rad/s}\ t + 1,04\ \mathrm{rad})\ \mathrm{m}\ \text{\punts{0,3}} \]
\end{solapartat}
